\documentclass[border=8pt]{standalone}
\usepackage{ctex}
\usepackage{amsmath}
\usepackage{tikz}
\begin{document}
\definecolor{mmRootFill}{RGB}{18,85,126}
\definecolor{mmLOneFill}{RGB}{234,148,136}
\begin{tikzpicture}[
  mmroot/.style={draw=black!75, line width=1.0pt, rounded corners=3.5pt, fill=mmRootFill, inner sep=5pt, font=\normalsize\bfseries},
  mml1/.style={draw=black!55, line width=0.9pt, rounded corners=2.5pt, fill=mmLOneFill, inner sep=4.5pt, font=\small\bfseries, align=left, text width=4.8cm},
  mml2/.style={draw=black!45, line width=0.8pt, rounded corners=2pt, fill=white, inner sep=4pt, font=\small, align=left, text width=6.2cm},
]
  %% 连线：根 → 一级 → 二级（节点填充不透明，盖住线头）
  \draw[black!55, line width=0.9pt] (4.35, 0.575) -- (4.8, 1.725);
  \draw[black!55, line width=0.9pt] (4.35, 0.575) -- (4.8, 0.575);
  \draw[black!55, line width=0.9pt] (4.35, 0.575) -- (4.8, -0.575);
  %% 节点
  \node[mmroot, anchor=east] at (4.35, 0.575) {2.4 曲线与方程};
  \node[mml1, anchor=west] at (4.6, 1.725) {曲线的方程与方程的曲线};
  \node[mml1, anchor=west] at (4.6, 0.575) {求动点的轨迹方程的一般步骤};
  \node[mml1, anchor=west] at (4.6, -0.575) {常见轨迹的判定结论};
\end{tikzpicture}
\end{document}
